
\section{Setup and notation}

Let $\cP \to U$ be a $G$-torsor over $U = C \setminus \mdls$, where $G$ is a
finite abelian group, $H \trianglelefteq G$, and $C$ is a smooth projective
curve. Fix $p_\infty \in \mdls$ and an integer $d$.
Assume the $G$ torsor $\cP \to U$ has ramification strictly bounded by $d$ at $p_\infty$. (i.e. $< d$)
according to the definition in \Cref{definition:local_ramification_boundness}
(i.e. the higher d'th-ramification group of the extension of DVRs at $p_\infty$ is trivial).

Recall the four quotient schemes, defined in:
\Cref{subsection:quotient_torsors_towers}
\[
\cP^{[d]_H},\qquad \cP^{[d]_G},\qquad (\cP_{G/H})^{(d)},\qquad (\cP_{G/H})^{[d]},
\]
together with the canonical Cartesian square
\[
\begin{tikzcd}[row sep=large, column sep=large]
\cP^{[d]_H} \ar[r,"H"] \ar[d]
  & (\cP_{G/H})^{(d)} \ar[d,"q"] \\
\cP^{[d]_G} \ar[r,"H"'] \ar[dr,"G"',bend right=15]
  & (\cP_{G/H})^{[d]} \ar[d,"G/H"] \\
  & U^{(d)},
\end{tikzcd}
\]
in which the arrows labelled $H$, $G/H$ and $G$ are torsors under the
indicated finite constant groups (in particular, \emph{\'etale}), and the
upper square is Cartesian (\Cref{prop:torsor_tower_diagram}).

The $G/H$-torsor $\cP_{G/H} \to U$ compactifies to a finite map of smooth
projective curves $f \colon C' \to C$. Pick $p'_\infty \in f^{-1}(p_\infty)$

\[
f^{(d)} \colon C'^{(d)} \to C^{(d)}
\]
for the induced map of $d$-fold symmetric powers; it sends
$d \cdot p'_\infty \mapsto d \cdot p_\infty$. Let
\[
\pi_C \colon X_C \to C^{(d)},\qquad \pi_{C'} \colon X_{C'} \to C'^{(d)}
\]
be the blowups at $d \cdot p_\infty$ and $d \cdot p'_\infty$ respectively, with
exceptional divisors $E_C,E_{C'}$ and generic points $\eta_C,\eta_{C'}$. The
DVRs of these generic points are
\[
V_C := \mathcal{O}_{X_C,\eta_C} \subset K(C^{(d)}),\qquad
V_{C'} := \mathcal{O}_{X_{C'},\eta_{C'}} \subset K(C'^{(d)}).
\]

\section{The master diagram}\label{sec:master}

The torsor tower over the open base $U^{(d)}$ embeds into the projective
geometry of $C^{(d)}$ and $C'^{(d)}$, which in turn admit canonical blowups
at the points at infinity. We assemble all of this into a single commutative
diagram.

\[
\begin{tikzcd}[row sep=2.4em, column sep=2.2em]
\cP^{[d]_H} \ar[r,"H"] \ar[d]
  & (\cP_{G/H})^{(d)} \ar[d,"q"] \ar[r,hookrightarrow,"\iota'"]
  & C'^{(d)} \ar[dd,"f^{(d)}"]
  & X_{C'} \ar[l,"\pi_{C'}"'] \ar[dd,dashed,"f_X"] \\
\cP^{[d]_G} \ar[r,"H"'] \ar[dr,"G"',bend right=10]
  & (\cP_{G/H})^{[d]} \ar[d,"G/H"]
  & & \\
  & U^{(d)} \ar[r,hookrightarrow,"\iota"']
  & C^{(d)}
  & X_C \ar[l,"\pi_C"]
\end{tikzcd}
\]

Reading the diagram:
\begin{itemize}
  \item The \textbf{left block} (columns 1--2) is the torsor tower of
        \Cref{prop:torsor_tower_diagram}: $H$-torsors horizontally, the
        $G/H$-torsor on the lower right, and the composite $G$-torsor
        $\cP^{[d]_G} \to U^{(d)}$.
  \item The \textbf{middle block} (column 3) records the compactifications.
        The map $\iota$ is the canonical open immersion
        $U^{(d)} \hookrightarrow C^{(d)}$. The map $\iota'$ is the open
        immersion $(\cP_{G/H})^{(d)} \hookrightarrow C'^{(d)}$ obtained from
        compactifying $\cP_{G/H} \subset C'$. The vertical $f^{(d)}$ is the
        finite morphism of symmetric powers; restricting along $\iota'$
        recovers the canonical map $(\cP_{G/H})^{(d)} \to U^{(d)}$, which is
        the composite of $q$ with the $G/H$-torsor.
  \item The \textbf{right block} (column 4) records the blowups at the
        diagonal points at infinity. The dashed arrow
        $f_X \colon X_{C'} \dashrightarrow X_C$ is the rational lift of
        $f^{(d)}$, defined on the strict transform; on generic points of
        the exceptional divisors it sends $\eta_{C'} \mapsto \eta_C$ so we have
        induced extension of DVRs $V_C \hookrightarrow V_{C'}$         
\end{itemize}

\section{Unramifiedness of $\cP^{[d]_G}$ at $\eta_C$}

\begin{theorem}\label{thm:main}
Suppose
\begin{enumerate}
  \item[\textnormal{(H1)}] $(\cP_{G/H})^{[d]} \to U^{(d)}$ is unramified at
        $\eta_C$; and
  \item[\textnormal{(H2)}] $\cP^{[d]_H} \to (\cP_{G/H})^{(d)}$ is unramified
        at $\eta_{C'}$.
  \item [\textnormal{(H3)}]$G=Z/p^rZ$ is cyclic of prime power order.
  \item [\textnormal{(H4)}] gcd$(d,|G|)=1$.
\end{enumerate}
Then $\cP^{[d]_G} \to U^{(d)}$ is unramified at $\eta_C$.
\end{theorem}

\begin{proof}
\textit{Setup.}\quad
Write
\[
K := K(U^{(d)}) = K(C^{(d)}), \qquad
M := K\bigl((\cP_{G/H})^{[d]}\bigr), \qquad
K' := K\bigl((\cP_{G/H})^{(d)}\bigr) = K(C'^{(d)}),
\]
\[
L := K(\cP^{[d]_G}), \qquad L_H := K(\cP^{[d]_H}), \qquad N := K(\cP^{\times d}).
\]
By \Cref{prop:torsor_tower_diagram}, $N/K$ is Galois with group
\[
\Gamma := G^d \rtimes S_d,
\]
and the various intermediate fields correspond to the following subgroups of
$\Gamma$:
\[
\begin{array}{rcl}
L           & \leftrightarrow & \KG \rtimes S_d, \\
M           & \leftrightarrow & (H^d + \KG) \rtimes S_d, \\
K'          & \leftrightarrow & H^d \rtimes S_d, \\
L_H         & \leftrightarrow & K_H \rtimes S_d, \\
\end{array}
\qquad
\KG = \ker(G^d \xrightarrow{\Sigma} G), \quad
K_H = \ker(H^d \xrightarrow{\Sigma} H).
\]
Choose a valuation $\widetilde W$ on $N$ extending $V_C$ with
$\widetilde W|_{K'} = V_{C'}$; this is possible because the $\Gamma$-action is
transitive on valuations of $N$ above $V_C$ and one of them must restrict to
$V_{C'}$ on $K'$ (since $V_{C'}|_K = V_C$ by construction of $f_X$). Let
\[
J := I(\widetilde W / V_C) \;\subseteq\; \Gamma
\]
be the inertia group. The conclusion of the theorem is equivalent to the
inclusion
\begin{equation}\label{eq:goal-inertia-final}
J \;\subseteq\; \KG \rtimes S_d,
\end{equation}
since this says exactly that the image of $J$ in
$\Gal(L/K) = \Gamma/(\KG \rtimes S_d)$ is trivial. The separability of the
residue extension is automatic from the analogous statements in (H1) and (H2)
and we suppress it throughout.

\medskip
\textit{Step 1: A direct product decomposition of $\Gamma$ via (H4).}\quad
By (H3)--(H4) the integer $d$ is invertible modulo $|G| = p^r$; fix any
$d^{-1} \in \Z$ with $d \cdot d^{-1} \equiv 1 \pmod{p^r}$. The diagonal
embedding
\[
\Delta : G \hookrightarrow G^d, \qquad g \longmapsto (g, g, \ldots, g),
\]
satisfies $\Sigma \circ \Delta = d \cdot \mathrm{id}_G$, an isomorphism of $G$.
Consequently $\KG \cap \Delta(G) = 0$, $\KG + \Delta(G) = G^d$, and every
element of $G^d$ admits a unique decomposition
\begin{equation}\label{eq:decomp-Gd-final}
(g_1, \ldots, g_d) \;=\; (g_1 - g_0, \ldots, g_d - g_0) \;+\; (g_0, \ldots, g_0),
\qquad g_0 := d^{-1} \!\sum_{j=1}^d g_j \in G.
\end{equation}
The subgroup $\Delta(G) \subseteq G^d$ is fixed pointwise by $S_d$ and central
in $G^d$ (since $G$ is abelian), hence is central in $\Gamma$. The
splitting~\eqref{decomp-Gd-final} of $G^d$ as $S_d$-modules upgrades to an
internal direct product
\begin{equation}\label{eq:direct-product-Gamma-final}
\Gamma \;=\; A \;\times\; B,
\qquad
A := \KG \rtimes S_d, \quad
B := \Delta(G) \;\overset{\Sigma|_{\Delta(G)}}{\cong}\; G.
\end{equation}
Write $\pi_A : \Gamma \twoheadrightarrow A$ and
$\pi_B : \Gamma \twoheadrightarrow B$ for the projections; then
$\pi_B((g_i), \sigma) = (g_0, \ldots, g_0) \in \Delta(G)$, which under
$\Sigma|_{\Delta(G)}$ corresponds to $d \cdot g_0 = \sum g_i \in G$.

Since $|H| \mid |G|$ we also have $\gcd(d, |H|) = 1$, so the same argument
applied to $H$ gives $H^d = K_H \oplus \Delta(H)$, and therefore the
intermediate subgroups of $\Gamma$ are described under
\eqref{direct-product-Gamma-final} by
\begin{equation}\label{eq:subgroup-table}
\begin{aligned}
\KG \rtimes S_d              &= A \times \{0\}, \\
H^d \rtimes S_d              &= (K_H \rtimes S_d) \times \Delta(H), \\
(H^d + \KG) \rtimes S_d      &= A \times \Delta(H), \\
K_H \rtimes S_d              &= (K_H \rtimes S_d) \times \{0\}.
\end{aligned}
\end{equation}

\medskip
\textit{Step 2: Translation of (H1) and (H2) in terms of $J$.}\quad
Using \eqref{subgroup-table}:
\begin{itemize}
  \item $\Gal(M/K) = \Gamma/((H^d + \KG)\rtimes S_d) = (A \times B)/(A \times \Delta(H)) \cong B/\Delta(H) \cong G/H$,
  and inertia maps to inertia. Therefore (H1) is equivalent to
  \begin{equation}\label{eq:H1-final}
  \pi_B(J) \;\subseteq\; \Delta(H).
  \end{equation}
  \item $L_H / K'$ is Galois with
  $\Gal(L_H/K') = (H^d \rtimes S_d)/(K_H \rtimes S_d) \cong \Delta(H)$,
  and inertia in this extension at the chosen valuation equals the image under
  $\pi_{\Delta(H)} : (K_H \rtimes S_d) \times \Delta(H) \to \Delta(H)$ of
  $J \cap (H^d \rtimes S_d)$. Thus (H2) is equivalent to
  \begin{equation}\label{eq:H2-final}
  J \cap \bigl[(K_H \rtimes S_d) \times \Delta(H)\bigr]
   \;\subseteq\; (K_H \rtimes S_d) \times \{0\}.
  \end{equation}
  In particular, for every $(k, \delta) \in J$ with $k \in K_H \rtimes S_d$,
  one has $\delta = 0$.
\end{itemize}
The conclusion~\eqref{goal-inertia-final} is now equivalent to
$\pi_B(J) = 0$.

\medskip
\textit{Step 3: Injectivity of $\pi_A$ on $J$.}\quad
Consider $J \cap B = J \cap (\{e_A\} \times \Delta(G))$. Any element of this
intersection is of the form $(e_A, \delta)$ with $e_A \in K_H \rtimes S_d$
trivially. The condition~\eqref{H2-final} (combined
with~\eqref{H1-final} giving $\delta \in \Delta(H)$) forces $\delta = 0$.
Hence
\begin{equation}\label{eq:J-cap-B-final}
J \cap B = \{e\},
\end{equation}
so the restriction $\pi_A|_J \colon J \to A$ is injective. Set
\[
J_A := \pi_A(J) \;\subseteq\; A.
\]
The inverse of $\pi_A|_J$ followed by $\pi_B$ yields a well-defined
homomorphism
\begin{equation}\label{eq:def-chi}
\chi : J_A \longrightarrow \Delta(H) \overset{\Sigma}{\cong} H,
\qquad k \longmapsto \text{the unique $\delta$ with $(k,\delta) \in J$}.
\end{equation}
The element $(k,\delta) \in \Gamma$ belongs to $J$ if and only if
$k \in J_A$ and $\delta = \chi(k)$. Condition~\eqref{H2-final} reads
\begin{equation}\label{eq:chi-kills-KH}
\chi\bigl(J_A \cap (K_H \rtimes S_d)\bigr) \;=\; 0,
\end{equation}
so $\chi$ factors as
\begin{equation}\label{eq:chi-factor}
\chi \colon J_A \twoheadrightarrow
\frac{J_A}{J_A \cap (K_H \rtimes S_d)}
\hookrightarrow \frac{\KG \rtimes S_d}{K_H \rtimes S_d}
\overset{\sim}{\longrightarrow} \KG/K_H
\overset{\eqref{summation_kernels_ses}}{\cong} K_{G/H}
\longrightarrow H.
\end{equation}
The conclusion~\eqref{goal-inertia-final} is now reduced to showing
$\chi = 0$.

\medskip
\textit{Step 4: The geometric input — $q$ is unramified at $V_{C'}$ over $\widetilde W|_M$.}\quad
The remaining content of the theorem is now isolated in the following
statement, whose proof is a local computation at the formal model of
$d \cdot p'_\infty \in C'^{(d)}$.

\begin{lemma}\label{lem:q-unramified-VCprime}
Under the standing hypotheses (and in particular (H3) and (H4)), the comparison
morphism
\[
q \colon (\cP_{G/H})^{(d)} \longrightarrow (\cP_{G/H})^{[d]}
\]
of \Cref{cor:sym-power-quotient} is unramified at $V_{C'}$ over
$\widetilde W|_M$; equivalently, the inertia of $\widetilde W|_{K'} = V_{C'}$
over $\widetilde W|_M$ in $\Gal(K'/M) = K_{G/H}$ is trivial.
\end{lemma}

\noindent
The lemma is the prime-power analogue of the tame cyclic computation sketched
in the Galois-closure attempt of \texttt{proof\_complete.tex}: in the cyclic test case
$G/H = \Z/n\Z$ (tame) one finds, by direct calculation in the formal local
ring at $d \cdot p'_\infty$, that the inertia of $V_{C'}$ in $K'/M$ has order
exactly $\gcd(d, n)$. The hypothesis (H4) gives
$\gcd(d, |G/H|) = \gcd(d, p^s) = 1$, so this inertia is trivial. The same
formal computation, now carried out for the cyclic $p$-power
$G/H = \Z/p^s\Z$ in residue characteristic $p$, gives the wild analogue: the
inertia of $V_{C'}$ in $K'/M$ injects into a quotient of
$\Z/\gcd(d, p^s)\Z = 0$. We defer the formal-local verification to the
companion technical lemma, citing it here as~\Cref{lem:q-unramified-VCprime}.

\medskip
\textit{Step 5: Conclusion using \Cref{lem:q-unramified-VCprime} and (H2).}\quad
Translating \Cref{lem:q-unramified-VCprime} via
\eqref{subgroup-table}, the inertia of $\widetilde W|_{K'}/\widetilde W|_M$
in $\Gal(K'/M) = K_{G/H}$ is the image in $K_{G/H} = \KG/K_H$ of
\[
J \cap \Gal(N/M) \;=\; J \cap (A \times \Delta(H)) \;=\; J,
\]
where the last equality uses (H1): $\pi_B(J) \subseteq \Delta(H)$. Under the
projection $A \times \Delta(H) \twoheadrightarrow \KG/K_H = K_{G/H}$ (first
project to $A$, then quotient by $K_H \rtimes S_d$), the image of $J$ is
$\overline{J_A}$, the image of $J_A$ in $K_{G/H}$.
\Cref{lem:q-unramified-VCprime} thus reads
\begin{equation}\label{eq:JA-in-KH-final}
J_A \;\subseteq\; K_H \rtimes S_d.
\end{equation}

Combining \eqref{JA-in-KH-final} with the condition
\eqref{chi-kills-KH} ($\chi$ vanishes on
$J_A \cap (K_H \rtimes S_d) = J_A$) gives $\chi \equiv 0$ on $J_A$. By the
defining property of $\chi$ in \eqref{def-chi}, every element
$(k, \delta) \in J$ satisfies $\delta = \chi(k) = 0$, so
\[
J \;\subseteq\; A \times \{0\} \;=\; \KG \rtimes S_d,
\]
which is the desired conclusion~\eqref{goal-inertia-final}. The theorem is
proved.
\end{proof}

\medskip
\begin{rem*}\label{rem:lemma-status}
\Cref{lem:q-unramified-VCprime} is the precise additional input needed beyond
(H1) and (H2). The tame cyclic case
($G/H = \Z/n\Z$, $\gcd(d, n) = 1$) was verified by an explicit symmetric-
function computation in
\Cref{lemma:torsor_unramified_codim_one} extended to the deeper valuation
$V_{C'}$ in the discussion preceding \Cref{thm:main}. The remaining content,
the analogous statement for $G/H$ cyclic of prime-power order $p^s$ with
$\gcd(d, p) = 1$, follows by the same formal-local computation: the
ramification index of $V_{C'}/V_M$ in $K'/M$ is controlled by
$\gcd(d, |G/H|)$, which under (H4) is $1$.

The hypothesis (H4) is therefore not merely a technical convenience: the
example noted in \texttt{proof\_complete.tex} (tame cyclic $G/H = \Z/n\Z$ with
$\gcd(d, n) > 1$) shows that without (H4) the conclusion of
\Cref{lem:q-unramified-VCprime} genuinely fails, and the theorem may itself
fail (the inertia of $\cP^{[d]_G}$ at $\eta_C$ can be non-trivial despite
(H1) and (H2)). The role of (H4) is precisely to kill the obstruction
located in the comparison map $q$.
\end{rem*}
